\begin{spacing}{1}
    \chapter*{Abstract}
\end{spacing}
\begin{wrapfigure}{r}{0.3\textwidth}
    \begin{center}
      \includegraphics[width=0.2\textwidth]{pics/question_mark.png}
    \end{center}
\end{wrapfigure}

\textbf{Task}

Modern companies operate complex IT infrastructures consisting of numerous servers,
services, and resources. Within large infrastructures, information about available
systems and their current state is often distributed across multiple platforms.
This makes it difficult for employees to maintain an overview and efficiently
analyze infrastructure-related data.

As part of this thesis, we present a specific solution for Primetals Technologies Austria GmbH.
The objective of the project is the development of a Global Infrastructure Overview (GIO) system,
which provides a centralized view over the company's IT infrastructure. Data from various sources
is collected, aggregated and visualized in a user-friendly manner.

\textbf{Implementation}

The implementation was done using a custom grafana plugin, which visualizes the collected data
on a world map and through detail views. The plugin allows the employees to filter and search for 
specific information. The data is collected from various grafana datasources, including MSSQL 
databases and enterprise APIs, through a centralized datasource layer before being provided to 
the visualization component.

\textbf{Result}

A system with the required functionalities was developed. The system is capable of visualizing
the collected data, so that the employees can efficiently analyse the current situation. 
The system has already been deployed at Primetals Technologies Austria GmbH, where it is 
actively used by employees to monitor the company's infrastructure. The system
can be extended with additional datasources and views.

\newpage
\begin{spacing}{1}
    \chapter*{Zusammenfassung}
\end{spacing}

\begin{wrapfigure}{r}{0.3\textwidth}
    \begin{center}
      \includegraphics[width=0.2\textwidth]{pics/question_mark.png}
    \end{center}
\end{wrapfigure}

\textbf{Aufgabenstellung}

Moderne Unternehmen verfügen über komplexe IT-Infrastrukturen, welche aus einer
Vielzahl an Servern, Diensten und weiteren Ressourcen bestehen. Durch die stetig
wachsende Anzahl an Systemen werden relevante Informationen häufig in
unterschiedlichen Plattformen und Datenquellen verwaltet. Dies erschwert es
Mitarbeitern, einen zentralen Überblick über die vorhandene Infrastruktur und
deren aktuellen Zustand zu erhalten.

Das Ziel dieser Arbeit ist die Entwicklung eines Global Infrastructure Overview
(GIO) Systems für die Primetals Technologies Austria GmbH. Das System soll Daten
aus unterschiedlichen Quellen zentral sammeln, aufbereiten und visualisieren,
um eine effizientere Analyse und Überwachung der Unternehmensinfrastruktur zu
ermöglichen.

\textbf{Realisierung}

Die Umsetzung erfolgte durch die Entwicklung eines zentralen Datasource-Layers
zur Aggregation und Verarbeitung von Infrastrukturinformationen aus
verschiedenen Quellen. Dazu werden Daten unter anderem aus MSSQL-Datenbanken und
Enterprise APIs gesammelt, vereinheitlicht und für die weitere Verarbeitung
bereitgestellt.

Für die Visualisierung wurde ein eigenes Grafana Plugin entwickelt. Dieses
ermöglicht die Darstellung der gesammelten Informationen auf einer Weltkarte
sowie in detaillierten Ansichten. Zusätzlich wurden Funktionen zur Filterung und
Suche implementiert, wodurch Mitarbeiter gezielt nach bestimmten
Infrastrukturinformationen suchen können.

\textbf{Ergebnis}

Es wurde ein System entwickelt, welches eine zentrale Übersicht über die
IT-Infrastruktur der Primetals Technologies Austria GmbH ermöglicht. Die
gesammelten Daten können mithilfe verschiedener Visualisierungen analysiert und
nach relevanten Informationen gefiltert werden.

Das System wurde erfolgreich in die bestehende Infrastruktur integriert und wird
von Mitarbeitern zur Überwachung und Analyse verwendet. Die Lösung kann durch
zusätzliche Datenquellen und weitere Visualisierungen erweitert werden.
